\section{Regressione non lineare nel caso scalare}

\subsection{Formulazione del problema e impossibilità della soluzione analitica}

Quando l'uscita di un modello matematico $y(t)$ viene descritta da un insieme di parametri non lineari, il problema della stima dei parametri diventa considerevolmente più complesso rispetto al caso lineare.

Per illustrare il problema, consideriamo come esempio il modello esponenziale scalare:
\begin{equation}
y(t) = g(t, \alpha) = e^{-\alpha t}
\end{equation}
dove $\alpha$ è un parametro non lineare da stimare a partire dalle misure sperimentali $z_{i}$ rilevate agli istanti $t_{i}$ ($i = 1, 2, \dots, N$).

Secondo il criterio dei minimi quadrati pesati, l'espressione (funzione di costo) da minimizzare rispetto al parametro $\alpha$ è la somma pesata dei residui quadrati (\textit{Weighted Residual Sum of Squares}, $WRSS$):
\begin{equation}
WRSS(\alpha) = \sum_{i=1}^{N} w_{i}(z_{i} - y_{i})^{2} = \sum_{i=1}^{N} w_{i}(z_{i} - e^{-\alpha t_{i}})^{2}
\end{equation}
dove $w_{i}$ rappresentano i pesi associati alle singole osservazioni.

Poiché $y(t)$ è una funzione non lineare nel parametro $\alpha$, non è possibile ricavare una soluzione analitica in forma chiusa per il problema di minimizzazione. Infatti, calcolando la derivata prima di $WRSS(\alpha)$ rispetto ad $\alpha$ e imponendo la condizione di stazionarietà uguale a zero, si ottiene:
\begin{equation}
\frac{d\,WRSS(\alpha)}{d\alpha} = \frac{d\left(\sum_{i=1}^{N} w_{i}(z_{i} - e^{-\alpha t_{i}})^{2}\right)}{d\alpha} = 0
\end{equation}
sviluppando la derivata:
\begin{equation}
-2 \sum_{i=1}^{N} w_{i} t_{i} e^{-\alpha t_{i}} (z_{i} - e^{-\alpha t_{i}}) = 0
\end{equation}
Tale equazione trascendente non consente di isolare algebricamente l'incognita e non fornisce quindi un'espressione esplicita per $\alpha$.

\subsection{Linearizzazione iterativa mediante il metodo di Gauss-Newton}

Per superare l'assenza di una soluzione analitica chiusa, una possibile strategia consiste nel realizzare una linearizzazione iterativa del modello attraverso il \textbf{metodo di Gauss-Newton}.

Si esegue innanzitutto l'espansione in serie di Taylor della funzione modello $y(t) = g(\alpha, t)$ intorno ad uno specifico valore iniziale $\alpha = \alpha^{0}$:
\begin{equation}
y(t) = g(\alpha, t) = g(\alpha^{0}, t) + \left.\frac{dg(\alpha, t)}{d\alpha}\right|_{\alpha=\alpha^{0}} (\alpha - \alpha^{0}) + \frac{1}{2}\left.\frac{d^{2}g(\alpha, t)}{d\alpha^{2}}\right|_{\alpha=\alpha^{0}} (\alpha - \alpha^{0})^{2} + \cdots
\end{equation}
Troncando lo sviluppo in serie al primo ordine (ovvero fermandosi alla derivata prima) e ipotizzando che la stima iniziale $\alpha^{0}$ sia nota, si ottiene l'approssimazione:
\begin{equation}
y(t) \approx \left.g(\alpha, t)\right|_{\alpha=\alpha^{0}} + \left.\frac{dg(\alpha, t)}{d\alpha}\right|_{\alpha=\alpha^{0}} (\alpha - \alpha^{0}) = e^{-\alpha^{0} t} - t e^{-\alpha^{0} t} (\alpha - \alpha^{0})
\end{equation}
Poiché $\alpha^{0}$ è noto, l'equazione risultante è ora lineare nel parametro $\alpha$.

Definendo lo scostamento parametrico:
\begin{equation}
\Delta\alpha = \alpha - \alpha^{0}
\end{equation}
per il generico istante di campionamento $t_{i}$ l'uscita del modello linearizzato assume la forma:
\begin{equation}
y_{i} \approx \left.g(\alpha, t)\right|_{\substack{\alpha=\alpha^{0} \\ t=t_{i}}} + \left.\frac{dg(\alpha, t)}{d\alpha}\right|_{\substack{\alpha=\alpha^{0} \\ t=t_{i}}} \Delta\alpha = e^{-\alpha^{0} t_{i}} - t_{i} e^{-\alpha^{0} t_{i}} \Delta\alpha
\end{equation}

Sostituendo questa espressione linearizzata all'interno della funzione di costo $WRSS(\alpha)$, si ricava:
\begin{align}
WRSS(\alpha) &= \sum_{i=1}^{N} w_{i}(z_{i} - y_{i})^{2} \nonumber \\
&\approx \sum_{i=1}^{N} w_{i}\left(z_{i} - \left.g(\alpha, t)\right|_{\substack{\alpha=\alpha^{0} \\ t=t_{i}}} - \left.\frac{dg(\alpha, t)}{d\alpha}\right|_{\substack{\alpha=\alpha^{0} \\ t=t_{i}}} \Delta\alpha\right)^{2} \nonumber \\
&= \sum_{i=1}^{N} w_{i}\left(z_{i} - e^{-\alpha^{0} t_{i}} + t_{i} e^{-\alpha^{0} t_{i}} \Delta\alpha\right)^{2}
\end{align}

Per convenienza notazionale, definiamo il residuo iniziale $\Delta z_{i}$ come la differenza tra il dato misurato $z_{i}$ e la predizione del modello calcolata nella stima iniziale $\alpha^{0}$:
\begin{equation}
\Delta z_{i} = z_{i} - y(\alpha^{0}, t_{i}) = z_{i} - e^{-\alpha^{0} t_{i}}
\end{equation}
Sostituendo $\Delta z_{i}$ nell'espressione approssimata di $WRSS$, si ottiene:
\begin{equation}
WRSS(\alpha) \approx \sum_{i=1}^{N} w_{i}\left(\underbrace{z_{i} - e^{-\alpha^{0} t_{i}}}_{\Delta z_{i}} + t_{i} e^{-\alpha^{0} t_{i}} \Delta\alpha\right)^{2} = \sum_{i=1}^{N} w_{i}\left(\Delta z_{i} + t_{i} e^{-\alpha^{0} t_{i}} \Delta\alpha\right)^{2}
\end{equation}
In questa formulazione, il modello dei residui è lineare nell'incremento $\Delta\alpha$ (e la funzione $WRSS$ ridotta al problema ai minimi quadrati lineare in $\Delta\alpha$ permette di calcolare analiticamente la correzione ottima $\Delta\hat{\alpha}$).

\subsection{Aggiornamento iterativo del parametro e minimi locali}

Minimizzando rispetto a $\Delta\alpha$, la stima dell'incremento ottimo $\Delta\hat{\alpha}$ è data da:
\begin{equation}
\Delta\hat{\alpha} = \frac{\sum_{i=1}^{N} w_{i} \Delta z_{i} \left.\frac{dg(\alpha, t)}{d\alpha}\right|_{\substack{\alpha=\alpha^{0} \\ t=t_{i}}}}{\sum_{i=1}^{N} w_{i} \left(\left.\frac{dg(\alpha, t)}{d\alpha}\right|_{\substack{\alpha=\alpha^{0} \\ t=t_{i}}}\right)^{2}} = \frac{-\sum_{i=1}^{N} w_{i} \Delta z_{i} t_{i} e^{-\alpha^{0} t_{i}}}{\sum_{i=1}^{N} w_{i} t_{i}^{2} e^{-2\alpha^{0} t_{i}}} = \frac{-\sum_{i=1}^{N} \frac{1}{\sigma_{i}^{2}} \Delta z_{i} t_{i} e^{-\alpha^{0} t_{i}}}{\sum_{i=1}^{N} \frac{1}{\sigma_{i}^{2}} t_{i}^{2} e^{-2\alpha^{0} t_{i}}}
\end{equation}
dove nell'ultimo passaggio si è esplicitata la scelta standard dei pesi $w_{i} = 1/\sigma_{i}^{2}$, inversamente proporzionali alla varianza $\sigma_{i}^{2}$ dell'errore di misura all'istante $t_{i}$.

Ricordando che $\Delta\alpha = \alpha - \alpha^{0}$, si ottiene una nuova stima aggiornata per il parametro $\alpha$:
\begin{equation}
\alpha^{1} = \alpha^{0} + \Delta\hat{\alpha}
\end{equation}
Il processo iterativo viene quindi ripetuto utilizzando $\alpha^{1}$ al posto di $\alpha^{0}$ nelle formule precedenti. Ad ogni passo dell'iterazione il valore di $WRSS$ si riduce progressivamente fino a raggiungere un punto di minimo, in corrispondenza del quale l'algoritmo iterativo si arresta.

È fondamentale evidenziare una differenza strutturale rispetto al caso lineare:
\begin{itemize}
    \item nella \textbf{regressione lineare} esiste un unico minimo per la funzione di costo $WRSS$;
    \item nella \textbf{regressione non lineare} l'andamento di $WRSS(\alpha)$ nel dominio del parametro può presentare più di un minimo locale (alcuni ben definiti o acuti, altri più graduali) e un solo minimo globale (adattato da Cobelli et al., 2001). Di conseguenza, la soluzione raggiunta dall'algoritmo iterativo è una soluzione approssimata che può convergere verso un minimo locale anziché al minimo globale, a seconda della scelta del punto iniziale $\alpha^{0}$.
\end{itemize}

\section{Estensione al caso vettoriale}

\subsection{Modello esponenziale a due parametri}

Estendiamo ora il procedimento al caso di un modello descritto da un vettore di due parametri, considerando la funzione:
\begin{equation}
y(t) = g(t, \mathbf{p}) = g(t, A, \alpha) = A e^{-\alpha t}
\end{equation}
Assumiamo di disporre di una stima iniziale del vettore dei parametri $\mathbf{p}^{0} = [A^{0}, \alpha^{0}]^{T}$. Sviluppando $g(t, A, \alpha)$ in serie di Taylor al primo ordine intorno a $(A^{0}, \alpha^{0})$, si ha:
\begin{equation}
g(t, A, \alpha) \cong g(t, A^{0}, \alpha^{0}) + \left.\frac{\partial g(t, A, \alpha)}{\partial \alpha}\right|_{\substack{A=A^{0} \\ \alpha=\alpha^{0}}} (\alpha - \alpha^{0}) + \left.\frac{\partial g(t, A, \alpha)}{\partial A}\right|_{\substack{A=A^{0} \\ \alpha=\alpha^{0}}} (A - A^{0})
\end{equation}
Calcolando il valore della funzione e delle sue derivate parziali nel punto iniziale $(A^{0}, \alpha^{0})$:
\begin{align}
g(t, A^{0}, \alpha^{0}) &= A^{0} e^{-\alpha^{0} t} \\
\left.\frac{\partial g(t, A, \alpha)}{\partial \alpha}\right|_{\substack{A=A^{0} \\ \alpha=\alpha^{0}}} &= -A^{0} t e^{-\alpha^{0} t} \\
\left.\frac{\partial g(t, A, \alpha)}{\partial A}\right|_{\substack{A=A^{0} \\ \alpha=\alpha^{0}}} &= e^{-\alpha^{0} t}
\end{align}
possiamo riscrivere l'approssimazione lineare come:
\begin{equation}
g(t, A, \alpha) \cong A^{0} e^{-\alpha^{0} t} - A^{0} t e^{-\alpha^{0} t} (\alpha - \alpha^{0}) + e^{-\alpha^{0} t} (A - A^{0})
\end{equation}

Consideriamo ora i dati sperimentali rumorosi acquisiti agli istanti $t_{i}$:
\begin{equation}
z_{i} = g(t_{i}, A, \alpha) + v_{i} = A e^{-\alpha t_{i}} + v_{i}
\end{equation}
dove $v_{i}$ rappresenta l'errore di misura. Utilizzando la linearizzazione precedente si ottiene:
\begin{equation}
z_{i} \cong A^{0} e^{-\alpha^{0} t_{i}} - A^{0} t_{i} e^{-\alpha^{0} t_{i}} (\alpha - \alpha^{0}) + e^{-\alpha^{0} t_{i}} (A - A^{0}) + v_{i}
\end{equation}
Definendo le variazioni incrementali:
\begin{equation}
\Delta z_{i} = z_{i} - A^{0} e^{-\alpha^{0} t_{i}}, \qquad \Delta\alpha = \alpha - \alpha^{0}, \qquad \Delta A = A - A^{0}
\end{equation}
l'equazione per la singola misura diventa:
\begin{equation}
\Delta z_{i} = -A^{0} t_{i} e^{-\alpha^{0} t_{i}} \Delta\alpha + e^{-\alpha^{0} t_{i}} \Delta A + v_{i}
\end{equation}

Considerando contemporaneamente tutte le $N$ misure ($i = 1, 2, \dots, N$), si ottiene la seguente formulazione matriciale:
\begin{equation}
\begin{bmatrix}
\Delta z_{1} \\
\Delta z_{2} \\
\vdots \\
\Delta z_{N-1} \\
\Delta z_{N}
\end{bmatrix}
=
\begin{bmatrix}
e^{-\alpha^{0} t_{1}} & -A^{0} t_{1} e^{-\alpha^{0} t_{1}} \\
e^{-\alpha^{0} t_{2}} & -A^{0} t_{2} e^{-\alpha^{0} t_{2}} \\
\vdots & \vdots \\
e^{-\alpha^{0} t_{N-1}} & -A^{0} t_{N-1} e^{-\alpha^{0} t_{N-1}} \\
e^{-\alpha^{0} t_{N}} & -A^{0} t_{N} e^{-\alpha^{0} t_{N}}
\end{bmatrix}
\begin{bmatrix}
\Delta A \\
\Delta\alpha
\end{bmatrix}
+
\begin{bmatrix}
v_{1} \\
v_{2} \\
\vdots \\
v_{N-1} \\
v_{N}
\end{bmatrix}
\end{equation}
In forma compatta, tale formulazione è lineare nel vettore degli incrementi parametrici $\Delta\mathbf{p} = [\Delta A, \Delta\alpha]^{T}$:
\begin{equation}
\Delta\mathbf{z} = \mathbf{G}\Delta\mathbf{p} + \mathbf{v}
\end{equation}
La soluzione ai minimi quadrati pesati (\textit{Weighted Least Squares}, WLS) è fornita da:
\begin{equation}
\Delta\hat{\mathbf{p}} = (\mathbf{G}^{T}\mathbf{\Sigma}_{v}^{-1}\mathbf{G})^{-1}\mathbf{G}^{T}\mathbf{\Sigma}_{v}^{-1}\Delta\mathbf{z}
\end{equation}
dove $\mathbf{\Sigma}_{v}$ è la matrice di covarianza del vettore di rumore $\mathbf{v}$. Una volta stimato $\Delta\hat{\mathbf{p}}$, si aggiorna la stima del vettore dei parametri:
\begin{equation}
\mathbf{p}^{1} = \mathbf{p}^{0} + \Delta\hat{\mathbf{p}}
\end{equation}
e si procede per iterazioni successive fino al soddisfacimento del criterio di arresto prefissato.

\subsection{Generalizzazione ad un generico modello non lineare a $M$ parametri}

Estendiamo ora la trattazione ad un generico modello non lineare descritto in forma vettoriale da:
\begin{equation}
\mathbf{z} = \mathbf{y} + \mathbf{v} = \mathbf{G}(\mathbf{p}) + \mathbf{v}
\end{equation}
dove:
\begin{itemize}
    \item $\mathbf{z} = [z_{1}, z_{2}, \dots, z_{N}]^{T}$ è il vettore delle $N$ misure;
    \item $\mathbf{y} = [y_{1}, y_{2}, \dots, y_{N}]^{T}$ è il vettore delle uscite ideali del modello;
    \item $\mathbf{p} = [p_{1}, p_{2}, \dots, p_{M}]^{T}$ è il vettore degli $M$ parametri incogniti;
    \item $\mathbf{G}(\mathbf{p}) = [g(t_{1}, \mathbf{p}), g(t_{2}, \mathbf{p}), \dots, g(t_{N}, \mathbf{p})]^{T}$ è il vettore delle predizioni non lineari;
    \item $\mathbf{v} = [v_{1}, v_{2}, \dots, v_{N}]^{T}$ è il vettore degli errori di misura con matrice di covarianza $\mathbf{\Sigma}_{v} = \operatorname{cov}(\mathbf{v})$.
\end{itemize}

La stima WLS dei parametri è il vettore $\hat{\mathbf{p}}_{WLS}$ che minimizza la funzione di costo quadratica pesata:
\begin{equation}
\hat{\mathbf{p}}_{WLS} = \arg\min_{\mathbf{p}} [\mathbf{z} - \mathbf{G}(\mathbf{p})]^{T}\mathbf{\Sigma}_{v}^{-1}[\mathbf{z} - \mathbf{G}(\mathbf{p})]
\end{equation}

Definiamo $\mathbf{p}^{0} = [p_{1}^{0}, p_{2}^{0}, \dots, p_{M}^{0}]^{T}$ la stima iniziale del vettore dei parametri $\mathbf{p}$. Linearizzando al primo ordine intorno a $\mathbf{p}^{0}$, la predizione del modello al generico istante $t_{i}$ è:
\begin{equation}
g(t_{i}, \mathbf{p}) \cong g(t_{i}, \mathbf{p}^{0}) +
\begin{bmatrix}
\dfrac{\partial g(t_{i}, \mathbf{p}^{0})}{\partial p_{1}} & \dfrac{\partial g(t_{i}, \mathbf{p}^{0})}{\partial p_{2}} & \cdots & \dfrac{\partial g(t_{i}, \mathbf{p}^{0})}{\partial p_{M}}
\end{bmatrix}
\begin{bmatrix}
p_{1} - p_{1}^{0} \\
p_{2} - p_{2}^{0} \\
\vdots \\
p_{M} - p_{M}^{0}
\end{bmatrix}
\end{equation}
Poiché le misure sono descritte da:
\begin{equation}
z_{i} = g(t_{i}, \mathbf{p}) + v_{i}, \qquad i = 1, 2, \dots, N
\end{equation}
sostituendo l'espressione linearizzata e raccogliendo tutte le $N$ equazioni in notazione vettoriale e matriciale, si ottiene:
\begin{equation}
\begin{bmatrix}
z_{1} - g(t_{1}, \mathbf{p}^{0}) \\
z_{2} - g(t_{2}, \mathbf{p}^{0}) \\
z_{3} - g(t_{3}, \mathbf{p}^{0}) \\
\vdots \\
z_{N} - g(t_{N}, \mathbf{p}^{0})
\end{bmatrix}
=
\begin{bmatrix}
\dfrac{\partial g(t_{1}, \mathbf{p}^{0})}{\partial p_{1}} & \dfrac{\partial g(t_{1}, \mathbf{p}^{0})}{\partial p_{2}} & \cdots & \dfrac{\partial g(t_{1}, \mathbf{p}^{0})}{\partial p_{M}} \\[2.2ex]
\dfrac{\partial g(t_{2}, \mathbf{p}^{0})}{\partial p_{1}} & \dfrac{\partial g(t_{2}, \mathbf{p}^{0})}{\partial p_{2}} & \cdots & \dfrac{\partial g(t_{2}, \mathbf{p}^{0})}{\partial p_{M}} \\[2.2ex]
\dfrac{\partial g(t_{3}, \mathbf{p}^{0})}{\partial p_{1}} & \dfrac{\partial g(t_{3}, \mathbf{p}^{0})}{\partial p_{2}} & \cdots & \dfrac{\partial g(t_{3}, \mathbf{p}^{0})}{\partial p_{M}} \\[1.5ex]
\vdots & \vdots & \ddots & \vdots \\[1.5ex]
\dfrac{\partial g(t_{N}, \mathbf{p}^{0})}{\partial p_{1}} & \dfrac{\partial g(t_{N}, \mathbf{p}^{0})}{\partial p_{2}} & \cdots & \dfrac{\partial g(t_{N}, \mathbf{p}^{0})}{\partial p_{M}}
\end{bmatrix}
\begin{bmatrix}
p_{1} - p_{1}^{0} \\
p_{2} - p_{2}^{0} \\
\vdots \\
p_{M} - p_{M}^{0}
\end{bmatrix}
+
\begin{bmatrix}
v_{1} \\
v_{2} \\
v_{3} \\
\vdots \\
v_{N}
\end{bmatrix}
\end{equation}

Si perviene così al sistema di equazioni lineari nei parametri incrementali $\Delta\mathbf{p} = \mathbf{p} - \mathbf{p}^{0}$:
\begin{equation}
\Delta\mathbf{z} = \mathbf{S}\Delta\mathbf{p} + \mathbf{v}
\end{equation}
dove $\mathbf{S}$ (di dimensione $N \times M$) è la \textbf{matrice di sensitività} (o matrice Jacobiana) valutata in $\mathbf{p}^{0}$, mentre $\Delta\mathbf{z} = \mathbf{z} - \mathbf{G}(\mathbf{p}^{0})$.

La soluzione WLS al passo corrente è data da:
\begin{equation}
\Delta\hat{\mathbf{p}} = (\mathbf{S}^{T}\mathbf{\Sigma}_{v}^{-1}\mathbf{S})^{-1}\mathbf{S}^{T}\mathbf{\Sigma}_{v}^{-1}\Delta\mathbf{z}
\end{equation}
e la nuova stima del vettore dei parametri si calcola come:
\begin{equation}
\mathbf{p}^{1} = \mathbf{p}^{0} + \Delta\hat{\mathbf{p}}
\end{equation}

Il processo iterativo può quindi ricominciare: il modello viene nuovamente linearizzato intorno a $\mathbf{p}^{1}$, si calcola una nuova matrice di sensitività e si ottiene una nuova stima $\mathbf{p}^{2}$. Si continua così fino a quando la funzione di costo smette di decrescere significativamente, ovvero quando la differenza tra due valori consecutivi di $WRSS(\mathbf{p})$ rientra all'interno di una tolleranza predefinita.

\section{Test per l'ordine del modello e principio di parsimonia}

\subsection{Criteri di informazione per la selezione dell'ordine}

Consideriamo il caso in cui si disponga di diversi modelli candidati e il problema consista nel selezionare quello che meglio approssima i dati sperimentali. Un esempio tipico è rappresentato dal modello multiesponenziale di una curva di decadimento:
\begin{equation}
y(t) = \sum_{i=1}^{n} A_{i} e^{-\lambda_{i} t}
\end{equation}
In molte applicazioni pratiche, l'ordine del modello --- che corrisponde al numero $n$ di componenti esponenziali --- non è noto a priori ed è pertanto necessario scegliere un valore ottimo per $n$.

Non è possibile definire l'ordine ottimo del modello basandosi esclusivamente sul valore minimo di $WRSS$ e sulla bianchezza dei residui. Infatti, all'aumentare dell'ordine del modello (ossia del numero di esponenziali), il valore di $WRSS$ decresce sistematicamente, ma contemporaneamente aumenta il numero $P$ di parametri da stimare. Di conseguenza, è necessario osservare congiuntamente sia il valore di $WRSS$ sia i gradi di libertà del sistema ($N - P$): bisogna cioè assicurarsi che la riduzione osservata in $WRSS$ rifletta un reale miglioramento nella descrizione del fenomeno e non sia dovuta semplicemente a una diminuzione dei gradi di libertà conseguente all'aumento del numero di parametri.

I due test più utilizzati nella pratica modellistica che implementano il \textbf{principio di parsimonia} (bilanciando bontà dell'adattamento e complessità del modello) sono:
\begin{enumerate}
    \item \textbf{Akaike Information Criterion (AIC)}:
    \begin{equation}
    AIC = WRSS + 2P
    \end{equation}
    \item \textbf{Schwartz Criterion (SC)}:
    \begin{equation}
    SC = WRSS + P \ln N
    \end{equation}
\end{enumerate}
dove:
\begin{itemize}
    \item $P \rightarrow$ numero totale dei parametri del modello;
    \item $N \rightarrow$ numero dei punti di misura.
\end{itemize}

\subsection{Esempio applicativo: scelta tra modello bi- e tri-esponenziale}

Consideriamo come esempio il confronto tra due modelli candidati --- un modello bi-esponenziale ($n = 2$, $P = 4$) e un modello tri-esponenziale ($n = 3$, $P = 6$) --- per descrivere un insieme di dati sperimentali di decadimento (adattato da Sparacino et al., 2000).

La Tabella~\ref{tab:parsimony} riporta i valori stimati per le ampiezze $A_{i}$ e le costanti di decadimento $\alpha_{i}$ (con i relativi indici di precisione percentuale tra parentesi), insieme ai corrispondenti valori di $WRSS$, $AIC$ e $SC$ ottenuti per i due modelli candidati.

\begin{table}[htbp]
\centering
\caption{Scelta dell'ordine del modello basata sul principio di parsimonia (adattato da Sparacino et al., 2000).}
\label{tab:parsimony}
{\footnotesize
\setlength{\tabcolsep}{3.2pt}
\begin{tabular}{@{}lccccccccc@{}}
\toprule
Modello & $A_{1}$ & $A_{2}$ & $A_{3}$ & $\alpha_{1}$ & $\alpha_{2}$ & $\alpha_{3}$ & $WRSS$ & $AIC$ & $SC$ \\
 & $(\mathrm{ml}^{-1})$ & $(\mathrm{ml}^{-1})$ & $(\mathrm{ml}^{-1})$ & $(\mathrm{min}^{-1})$ & $(\mathrm{min}^{-1})$ & $(\mathrm{min}^{-1})$ & & & \\
\midrule
2 Esponenziali & $1.79 \times 10^{-4}$ & $0.90 \times 10^{-4}$ & --- & $0.141$ & $0.025$ & --- & $61$ & $69$ & $75$ \\
 & $(7.2)$ & $(4.8)$ & & $(9.2)$ & $(2.5)$ & & & & \\
\addlinespace[4pt]
3 Esponenziali & $1.29 \times 10^{-4}$ & $0.88 \times 10^{-4}$ & $0.81 \times 10^{-4}$ & $0.253$ & $0.085$ & $0.024$ & $40$ & $52$ & $60$ \\
 & $(39)$ & $(66)$ & $(9)$ & $(49)$ & $(40)$ & $(3)$ & & & \\
\bottomrule
\end{tabular}
}
\end{table}

Confrontando gli indici riportati nella Tabella~\ref{tab:parsimony}, il modello più parsimonioso risulta essere quello caratterizzato da \textbf{3 esponenziali}, poiché presenta i valori più bassi sia per il criterio di Akaike ($AIC = 52$ contro $69$) sia per il criterio di Schwartz ($SC = 60$ contro $75$).